Flowering time in the plant species \textit{Arabidopsis thaliana} is a classic model for studying the evolution of plasticity and \gxe\  \citep{mitchell2006genetic, maple2024flowering, leventhal2025planting}.
As would be expected for an annual plant with a large geographic range, \textit{A. thaliana} ecotypes exhibit a range of life history strategies and these strategies can often be connected to environmental or geographic variation, contributing to \gxe\ \citep{stinchcombe2004latitudinal, mitchell2006genetic, hancock2011adaptation,martinez2020functional}.
%In addition, many QTLs for flowering time are environment specific \citep{li2014genetic, li2010association}.
Vernalization, the cold temperatures experienced by overwintering plants, is an important cue for flowering time and \textit{A. thaliana} ecotypes differ in their sensitivity to vernalization \citep{shindo2006variation, sung2005remembering}.
%There is evidence for selection on one of the genes important for vernalization response, FRIGIDA \citep{Toomajian2006nonparametric}, however FRIGIDA variation alone cannot explain observed variation in response to vernalization \citep{wilczek2010genetic}.
Despite the extensive study of this trait, we lack a clear link between flowering time plasticity variation and long-term selection, limiting our understanding of how local adaptation has shaped variation in this system.

Big picture
Traits are shaped by both genetics and the environment, so determining the evolutionary forces shaping trait variation requires not only identifying genetic variation for traits but also determining how this variation differs across environments.
Genetic variation for environmental response (or plasticity) can contribute to genotype by environment interactions (\gxe) in determining phenotypes \citep{falconer1952problem, bradshaw1965evolutionary}.
\gxe\ has long been of interest to evolutionary biologists because it could represent adaptive phenotypic plasticity  \citep{schlichting1986evolution,scheiner1993genetics}.


While adaptation through the evolution of plastic responses could increase fitness in heterogeneous environments \citep{van1991evolution, baythavong2011linking, anderson2021selection}, the evolution of plasticity could be limited by costs or other factors \citep{dewitt1998costs, callahan2005plasticity, roberts2023weaker}.
Alternatively, \gxe\ could result from different selection pressures across environments that constrain trait expression in some environments but allow trait expression to drift in environments where selection is relaxed \citep{paaby2014cryptic}.
Distinguishing between these scenarios has been challenging.

Different conceptions of plasticity have contributed to different methods for investigating the evolutionary forces shaping \gxe.
First, plasticity could be measured using reaction norms: functions of trait value expression across environments \citep{freeman1973statistical,scheiner1993genetics, morrissey2016variation, arnold2019analyse, de2025partitioning}.
Reaction norms have been studied in a number of contexts and been used to test for selection on plasticity \citep{valdes2023spring, wagner2018plasticity}.
One limitation of reaction norms is often they require doing 'statistics on statistics', losing variance when estimating reaction norms, although the random regression approach alleviates this issue \citep{arnold2019analyse}.
Second, trait values measured across environments can be thought of as correlated traits, and methods developed for studying multivariate trait evolution can be applied to plasticity.
Studies of plasticity as a multivariate trait have found that the genetic basis of traits differs across environments \citep{yan2021genetic}, identified specific types of weather contributing to phenological \gxe\ \citep{MacQueen_2025}, and measured selection for plasticity\citep {anderson2021selection}.